\section{The Dilution Crisis}\label{sec:dilution}

Scientific output has been growing faster than the community's
capacity to review it~\cite{publishing-strain2024}, and paper mills
industrialized fabrication well before language
models~\cite{paper-mills2021}. What changed is the cost curve.
\emph{Obscurity-Is-Dead} measured how language models compress the
effort gap for reverse-engineering proprietary
devices~\cite{obscurity-is-dead}; the same compression applies to
producing a plausible paper. GPT-fabricated papers are already
indexed and ranked by scholarly search
engines~\cite{gpt-fabricated2024}.

The demand side is the incentive system. Publish-or-perish pays for
counted output, and hypercompetition selects for whatever maximizes
the count~\cite{perverse-incentives2017}. Language models did not
create that pressure; they gave it a zero-cost supply channel.
Dilution is misaligned incentives meeting free volume.

The consequence is that finding good sources is now the bottleneck of
research. Retrieval trusts proxies: indexed, cited, fluent. Generation
has learned to imitate all three. Language models fabricate
plausible references and make citation errors at documented
rates~\cite{fabricated-citations2023, ai-hallucination2023}; once such
a reference is indexed, nothing in today's retrieval chain stops it
from being cited onward, and models trained on the polluted pool
degrade further~\cite{model-collapse2024}, so dilution compounds. The
parallelism that lets one operator advance several strands equally
industrializes pollution (Section~\ref{sec:discussion});
per-activity accounting is what tells the two apart.
The diluted literature is, in our words, a data swamp one level up,
and graph-shaped metadata harvesting is its deployed
remedy~\cite{unhide-helmholtz-kg}.

When text is cheap, the scarce good is an original idea whose support
is checkable. We treat originality as a provenance property: a novelty
claim carries its recorded search horizon, that is, what was searched,
where, and when. Concretely, the tooling's registry sweeps (Crossref,
OpenAlex, arXiv, DataCite) only print candidates; the agent logs a
sweep as an activity, and at least the novelty sweep behind this paper
records its query terms. So ``no prior work found'' stops being an unfalsifiable
assertion and becomes a bounded one: no prior work found \emph{within
this recorded horizon}, which a sceptical reader can re-run, widen,
and possibly refute. Every claim must then be supported, or refuted,
by facts.

aiprov's answer makes claims first-class. Each assertion carries
\texttt{prov:wasGeneratedBy}, \texttt{prov:wasAttributedTo}, and a
ladder rung that separates ``the reference exists''
(machine-checkable) from ``the content supports the claim''
(machine-judged at most) from ``a human checked'' (human-only).
Refutation is declared but not yet recorded: \texttt{aiprov:contradicts}
exists in the vocabulary for a claim refuted by facts, but no command
writes it, no validation rule or rung uses it, and the graphs reported
here hold no such triple; the review runs keep ``verified, not
promoted'' only. This is
supply-side transparency. Artefacts carry a pedigree that fabrication
cannot cheaply forge: DOIs that resolve, content hashes, commits, and
registry-resolved identities. Filtering the pool can then shift from
fluency heuristics to verifiable metadata.
